\noindent
The earlier papers of this series asked how a relation \emph{begins}; this one asks how it does not \emph{end}---under what conditions a generative bond reproduces itself across time. Its problem is that self-continuation, taken alone, is morally indifferent: capital, the rose, and the parasitic bond all continue themselves no less than love does. Eglash's recursive return of value is therefore necessary but not sufficient. The paper supplies the missing temporal dimension through an axiom of sustainability cast in the language of \emph{geometric phase}---the vicious cycle a circle of zero holonomy, the good cycle a spiral that returns to its root while sublimating---and then asks how, through signs, such a spiral might be fostered. Its constructive thesis is that the \emph{appropriate poem}---the non-possessive interpretive cycle that generates real value---is the form of the good circle, and that generative justice, in the intimate domain, is for this reason necessarily poetic.
